%% RESUMO EM LINGUA ESTRANGEIRA
\begin{abstract}
This dummy dissertation outlines a research project on the use of artificial intelligence in the triage of submissions received by public ombudsman offices. The general objective is to analyse, in an exploratory way, how natural language processing techniques can support the classification and routing of citizen demands while preserving human oversight and the traceability of decisions. The theoretical framework brings together three recent studies on petition classification, complaint routing and the orchestration of agents with a human in the loop. The planned procedure is a structured literature review in open bibliographic sources, reading abstracts and comparing the institutional context and the task addressed by each study. Since it is only an outline, the work presents no results of its own and points to the expansion of the review and the definition of selection criteria as next steps.
\end{abstract}
